\documentclass[10pt,a4paper]{amsart}
\usepackage{amsmath,amssymb,amsthm,graphicx}
\usepackage{verbatim}
\newenvironment{reference}{\comment}{\endcomment}
%\newenvironment{reference}{}{}
\newenvironment{slogan}{\comment}{\endcomment}
\newenvironment{history}{\comment}{\endcomment}

% For commutative diagrams we use Xy-pic
\usepackage[all]{xy}

% We use 2cell for 2-commutative diagrams.
\xyoption{2cell}
% Theorem environments.
%
\theoremstyle{plain}
\newtheorem{theorem}[subsection]{Theorem}
\newtheorem{proposition}[subsection]{Proposition}
\newtheorem{lemma}[subsection]{Lemma}

\theoremstyle{definition}
\newtheorem{definition}[subsection]{Definition}
\newtheorem{example}[subsection]{Example}
\newtheorem{exercise}[subsection]{Exercise}
\newtheorem{situation}[subsection]{Situation}

\theoremstyle{remark}
\newtheorem{remark}[subsection]{Remark}
\newtheorem{remarks}[subsection]{Remarks}

\numberwithin{equation}{subsection}

% Macros
%
\def\lim{\mathop{\mathrm{lim}}\nolimits}
\def\colim{\mathop{\mathrm{colim}}\nolimits}
\def\Spec{\mathop{\mathrm{Spec}}}
\def\Hom{\mathop{\mathrm{Hom}}\nolimits}
\def\Ext{\mathop{\mathrm{Ext}}\nolimits}
\def\SheafHom{\mathop{\mathcal{H}\!\mathit{om}}\nolimits}
\def\SheafExt{\mathop{\mathcal{E}\!\mathit{xt}}\nolimits}
\def\Sch{\mathit{Sch}}
\def\Mor{\mathop{\mathrm{Mor}}\nolimits}
\def\Ob{\mathop{\mathrm{Ob}}\nolimits}
\def\Sh{\mathop{\mathit{Sh}}\nolimits}
\def\NL{\mathop{N\!L}\nolimits}
\def\CH{\mathop{\mathrm{CH}}\nolimits}
\def\proetale{{pro\text{-}\acute{e}tale}}
\def\etale{{\acute{e}tale}}
\def\QCoh{\mathit{QCoh}}
\def\Ker{\mathop{\mathrm{Ker}}}
\def\Im{\mathop{\mathrm{Im}}}
\def\Coker{\mathop{\mathrm{Coker}}}
\def\Coim{\mathop{\mathrm{Coim}}}

% Boxtimes
%
\DeclareMathSymbol{\boxtimes}{\mathbin}{AMSa}{"02}

%
% Macros for moduli stacks/spaces
%
\def\QCohstack{\mathcal{QC}\!\mathit{oh}}
\def\Cohstack{\mathcal{C}\!\mathit{oh}}
\def\Spacesstack{\mathcal{S}\!\mathit{paces}}
\def\Quotfunctor{\mathrm{Quot}}
\def\Hilbfunctor{\mathrm{Hilb}}
\def\Curvesstack{\mathcal{C}\!\mathit{urves}}
\def\Polarizedstack{\mathcal{P}\!\mathit{olarized}}
\def\Complexesstack{\mathcal{C}\!\mathit{omplexes}}
% \Pic is the operator that assigns to X its picard group, usage \Pic(X)
% \Picardstack_{X/B} denotes the Picard stack of X over B
% \Picardfunctor_{X/B} denotes the Picard functor of X over B
\def\Pic{\mathop{\mathrm{Pic}}\nolimits}
\def\Picardstack{\mathcal{P}\!\mathit{ic}}
\def\Picardfunctor{\mathrm{Pic}}
\def\Deformationcategory{\mathcal{D}\!\mathit{ef}}
\usepackage{cleveref}
\crefname{originalreference}{original reference}{original references}
\Crefname{originalreference}{Original reference}{Original references}
\makeatletter
\newlabel{lemma-duality-dim-1}{{0BS2}{source}{}{sourcepointer-1}{}}
\newlabel{lemma-duality-dim-1@cref}{{[originalreference][0][]0BS2}{[1][0][]source}{}{}{}}
\makeatother
\crefname{originalreference}{original reference}{original references}
\Crefname{originalreference}{Original reference}{Original references}
\makeatletter
\newlabel{lemma-rr}{{0BS6}{source}{}{sourcepointer-2}{}}
\newlabel{lemma-rr@cref}{{[originalreference][0][]0BS6}{[1][0][]source}{}{}{}}
\makeatother
\crefname{originalreference}{original reference}{original references}
\Crefname{originalreference}{Original reference}{Original references}
\makeatletter
\newlabel{lemma-uniformizer-works}{{0C1E}{source}{}{sourcepointer-3}{}}
\newlabel{lemma-uniformizer-works@cref}{{[originalreference][0][]0C1E}{[1][0][]source}{}{}{}}
\makeatother
\crefname{originallemma}{original lemma}{original lemmas}
\Crefname{originallemma}{Original lemma}{Original lemmas}
\makeatletter
\newlabel{varieties-lemma-check-invertible-sheaf-trivial}{{0B40}{source}{}{sourcepointer-4}{}}
\newlabel{varieties-lemma-check-invertible-sheaf-trivial@cref}{{[originallemma][0][]0B40}{[1][0][]source}{}{}{}}
\makeatother
\title{Vector fields on smooth proper curves: vanishing in genus at least 2, fixed-point counts in genus 0 and 1}
\author{The Stacks Project Authors}
\date{stacks-project:curves.tex:local-snapshot-20261012; GNU FDL 1.2 or later}
\begin{document}
\maketitle
\noindent\textit{Editorial note.} All delivered mathematical text is the author's own, copied from the pinned Stacks Project chapter file curves.tex; no statement or proof step was rewritten, reordered or removed.\par\medskip
\noindent\textit{Editorial note.} The counted proof is the author's complete argument: the derivation identification, the degree computations in all three genera, and the uniformizer-based fixed-point translation are all retained verbatim.\par\medskip
\noindent\textit{Editorial note.} The label lemma-smooth-vector-fields names the counted statement, and the proof environment immediately following it in the source is delivered verbatim as its proof.\par\medskip
\noindent\textit{Editorial note.} The counted statement is carried by the author's own lemma environment in the chapter source.\par\medskip
\noindent\textit{Allowed standard prerequisites (editorial).}\begin{itemize}
\item Derivations as twisted differentials. Same chapter, Lemmas lemma-duality-dim-1 and lemma-uniformizer-works. Identifies Der\_k(O\_X, O\_X) with Hom(omega\_X, O\_X) via D = theta composed with d, and makes dz a basis of Omega at a closed point for a uniformizer z.
\item Riemann-Roch degree formula. Same chapter, Lemma lemma-rr. Gives deg(omega\_X) = 2g - 2, the input to the vanishing and vanishing-locus computations.
\item Global sections of invertible sheaves on curves. Varieties, Lemma varieties-lemma-check-invertible-sheaf-trivial. Forces theta = 0 when Hom(omega\_X, O\_X) = H\textasciicircum{}0(omega\_X\textasciicircum{}tensor -1) has negative degree, and controls the vanishing divisor in the genus 0 and 1 cases.
\end{itemize}\medskip
\ifdefined\hypertarget\hypertarget{sourcepointer-1}{}\fi
\noindent\textit{Source reference.} Original reference 0BS2: sources/stacks-curves/pinned-original-source.tex, source lines 796--796. Prerequisite pointer invoked at proof line(s) 6897; it names a standard fact of the same chapter that the counted argument uses, and every step of the counted proof built on it is retained verbatim in the delivered body.\par\medskip
\ifdefined\hypertarget\hypertarget{sourcepointer-2}{}\fi
\noindent\textit{Source reference.} Original reference 0BS6: sources/stacks-curves/pinned-original-source.tex, source lines 1218--1218. Prerequisite pointer invoked at proof line(s) 6899; it names a standard fact of the same chapter that the counted argument uses, and every step of the counted proof built on it is retained verbatim in the delivered body.\par\medskip
\ifdefined\hypertarget\hypertarget{sourcepointer-3}{}\fi
\noindent\textit{Source reference.} Original reference 0C1E: sources/stacks-curves/pinned-original-source.tex, source lines 2629--2629. Prerequisite pointer invoked at proof line(s) 6911; it names a standard fact of the same chapter that the counted argument uses, and every step of the counted proof built on it is retained verbatim in the delivered body.\par\medskip
\ifdefined\hypertarget\hypertarget{sourcepointer-4}{}\fi
\noindent\textit{Source reference.} Original lemma 0B40: sources/stacks-chow/tags.txt, source lines 13983--13983. Prerequisite pointer invoked at proof line(s) 6902; it names a standard fact of a sibling chapter that the counted argument uses, and every step of the counted proof built on it is retained verbatim in the delivered body.\par\medskip
\begin{lemma}
\label{lemma-smooth-vector-fields}
Let $k$ be an algebraically closed field.
Let $X$ be a smooth, proper, connected curve over $k$.
Let $g$ be the genus of $X$.
\begin{enumerate}
\item If $g \geq 2$, then $\text{Der}_k(\mathcal{O}_X, \mathcal{O}_X)$
is zero,
\item if $g = 1$ and $D \in \text{Der}_k(\mathcal{O}_X, \mathcal{O}_X)$
is nonzero, then $D$ does not fix any closed point of $X$, and
\item if $g = 0$ and $D \in \text{Der}_k(\mathcal{O}_X, \mathcal{O}_X)$
is nonzero, then $D$ fixes at most $2$ closed points of $X$.
\end{enumerate}
\end{lemma}

\begin{proof}
Recall that we have a universal $k$-derivation
$d : \mathcal{O}_X \to \Omega_{X/k}$ and hence $D = \theta \circ d$
for some $\mathcal{O}_X$-linear map $\theta : \Omega_{X/k} \to \mathcal{O}_X$.
Recall that $\Omega_{X/k} \cong \omega_X$, see
Lemma \ref{lemma-duality-dim-1}.
By Riemann-Roch we have $\deg(\omega_X) = 2g - 2$
(Lemma \ref{lemma-rr}).
Thus we see that $\theta$ is forced to be zero
if $g > 1$ by Varieties, Lemma
\ref{varieties-lemma-check-invertible-sheaf-trivial}.
This proves part (1).
If $g = 1$, then a nonzero $\theta$ does not vanish anywhere and if
$g = 0$, then a nonzero $\theta$ vanishes in a divisor of degree $2$.
Thus parts (2) and (3) follow if we show that
vanishing of $\theta$ at a closed point $x \in X$ is
equivalent to the statement that $D$ fixes $x$ (as defined above).
Let $z \in \mathcal{O}_{X, x}$ be a uniformizer.
Then $dz$ is a basis element for $\Omega_{X, x}$, see
Lemma \ref{lemma-uniformizer-works}.
Since $D(z) = \theta(dz)$ we conclude.
\end{proof}

\end{document}
